\subsection{Single-asset generator validation}

We first characterized the SPY series and the HMM-WJ calibrated to its
training window (Fig.~\ref{fig:empirical_motivation};
Table~\ref{tab:descriptive_stats}). The observed series had heavy tails
and persistent absolute-return autocorrelation in both the training
and holdout windows. Individual raw-return correlations were small,
although the Ljung--Box test detected joint serial dependence
(Table~\ref{tab:descriptive_stats}). Under the counted transition
matrix, every state had an expected residence time of only $1$--$2$
steps, and the implied absolute-return autocorrelation decayed to near
zero by lag $5$ (Fig.~\ref{fig:model_internals}; Online
Appendix~\ref{sec:supp-generator}). Jump episodes extended visits to
the tail states and so provided a mechanism for more persistent
volatility (Algorithm~\ref{alg:hmmwj}). Calibration against the
$2014$--$2024$ history selected mean episode duration $\lambda=90$ and
jump probability $\epsilon=10^{-4}$ (Eq.~\eqref{eq:grid_search};
Table~\ref{tab:hyperparams}). Because $\epsilon$ mainly sets how many
paths contain an episode and the objective scored only jump-active
paths, the objective depended only weakly on $\epsilon$, and the search
selected its smallest grid value (Online
Appendix~\ref{sec:supp-gen-tuning}). The duration $\lambda$ was
likewise a fitted parameter of the objective, not an estimate of how
long observed high-volatility periods last.

% Main table tab:model_comparison is collected after the references.

We compared the calibrated HMM-WJ with seven alternative generators on
the training window, using $1{,}000$ paths of $2{,}766$ trading days
(Table~\ref{tab:model_comparison}; Fig.~\ref{fig:model_comparison}).
Three alternatives drew independent returns: bootstrap resampling of
the training returns with replacement, and Gaussian and Laplace
distributions fitted to them. The other four were HMM-NJ, a GARCH(1,1)
model with Gaussian innovations, the GRU benchmark~\cite{choEtAl2014gru}, and the HSMM
baseline adapted from Bulla and Bulla~\cite{bullaBulla2006}. The HMM-NJ and HMM-WJ models had KS pass rates within about one
percentage point of the fitted Laplace baseline but came much closer
to the observed kurtosis, which Laplace underestimated at $3.0$
against $7.7$. Both HMM variants had lower
Wasserstein-1 distances than Laplace, whereas Laplace had the lowest
Hellinger distance among the fitted generators, so the ranking
depended on the metric. Bootstrap resampling passed every KS test
because it reused the observed returns, while Gaussian sampling passed
neither distribution test. The GARCH(1,1) model had the lowest mean
absolute error between the observed and ensemble-average
absolute-return autocorrelation functions (ACF-MAE) but rarely passed
the distribution tests, and the GRU showed the same tradeoff. The HSMM
baseline passed only $82\%$ of KS tests and underestimated kurtosis by
more than one third, and its ACF-MAE equaled the HMM-NJ value despite
its flexible state durations. A separate ablation within HMM-WJ
isolated the emission distribution: Student-$t$ emissions produced
kurtosis of $7.5$, close to the observed $7.7$, whereas Normal
($\nu=\infty$) emissions underestimated it by about $29\%$
(Table~\ref{tab:emission_hsmm}). Both HMM variants also reproduced the
observed skewness of $-0.75$ without a skewness parameter, because the
state means and scales estimated from the asymmetric series carried
its asymmetry into the mixture (Online
Appendix~\ref{sec:supp-generator};
Table~\ref{tab:supp-generator-quantities}).

Jump episodes traded a small loss of marginal fit for better temporal
fit in training (Table~\ref{tab:model_comparison};
Fig.~\ref{fig:model_comparison}b). About $24\%$ of the
training-length HMM-WJ paths contained an episode, and paths without
one were equivalent to HMM-NJ (Fig.~\ref{fig:statistical_validation}).
Without episodes, the counted chain mixed so quickly that HMM-NJ had
an ACF-MAE close to that of the three independent-draw generators.
Across all HMM-WJ paths, ACF-MAE fell from $0.059$ for HMM-NJ to
$0.053$, with a small reduction in KS and AD pass rates. The episode
share, the fraction of forced steps, and the autocorrelation of
jump-active paths agreed with the calculations in Online
Appendix~\ref{sec:supp-generator}
(Tables~\ref{tab:supp-generator-quantities} and
\ref{tab:supp-acf-prediction}). Because the objective scored only
jump-active paths, the one-quarter share was not a calibration target;
it followed from $\epsilon$ and the path length. To examine how that share traded marginal against temporal fit, we
split $4{,}000$ simulated HMM-WJ paths into jump-active and no-jump
groups and evaluated mixtures of the two, weighted by a jump-active
fraction from zero to one (Online Appendix~\ref{sec:supp-gen-tuning}). Increasing that fraction lowered autocorrelation
error but also lowered KS and AD pass rates, while the Hill upper-tail
index remained stable. Kurtosis fell from above the observed value at
a fraction of zero to below it at one. Forced steps drew repeatedly
from the existing tail states without extending the extreme tail, so
they raised the variance more than the fourth moment (Online
Appendix~\ref{sec:supp-gen-tuning}). At the realized share of about one quarter, the mean path kurtosis of
$7.5$ lay slightly below the observed $7.7$
(Table~\ref{tab:model_comparison}).

To test whether the HMM results depended on the state partition, we
refitted both HMM variants for $N$ from $30$ to $200$ states
(Table~\ref{tab:sensitivity_n}). Each refit kept five states in each
tail set and re-ran the $(\epsilon,\lambda)$ search. The HMM-NJ pass
rates changed by less than one percentage point. As $N$ increased,
HMM-WJ autocorrelation error fell but so did its AD pass rate. The
tail sets that receive forced steps then covered a narrower quantile
band. Each forced step therefore drew a more extreme return, which
lowered autocorrelation error, while each tail state's emission
estimates rested on fewer observations, which lowered the AD pass rate
(Online Appendix~\ref{sec:supp-sensitivity}). The primary analysis used
$N=100$, which kept the HMM-WJ KS pass rate above $98\%$ and its AD
pass rate within a few percentage points of its best value while
capturing part of the autocorrelation improvement. The sweep used the
$2{,}766$-day SPY history alone. It therefore did not establish a state
count for the other assets, which also used $N=100$, or for shorter
histories, which would leave fewer observations per state.

To test whether the distributional fit held on data not used for
fitting, we compared $249$-day paths from the frozen models with the
$249$ held-out SPY returns from $2025$
(Table~\ref{tab:model_comparison};
Fig.~\ref{fig:statistical_validation}). On $249$ observations the
tests had less power, so the pass rates of the weaker generators rose
sharply and the differences among generators narrowed. The HMM-NJ and
HMM-WJ models retained KS pass rates of $96.6\%$ and $95.4\%$, with
kurtosis close to the observed value. At $\epsilon=10^{-4}$, few
$249$-day paths contained an episode (Online
Appendix~\ref{sec:supp-gen-renewal}), so this comparison mainly tested
the no-jump dynamics that both variants share. The fitted Laplace
baseline matched or exceeded both HMM variants on the KS and AD pass
rates and had equal or smaller Wasserstein-1 and Hellinger distances.
Its excess kurtosis of $2.8$, however, remained well below the observed
$6.9$, whereas both HMM variants reached $6.2$
(Table~\ref{tab:model_comparison}). To check whether the single-asset
fit carried over to other stocks, we repeated the fitting and
calibration of HMM-WJ for NVDA, JNJ, and JPM on the same training and
holdout windows (Table~\ref{tab:cross_asset}; Online
Appendix~\ref{sec:supp-cross-asset}). Training KS pass rates were near $99\%$
for all three, but holdout rates ranged from $74.0\%$ for JNJ to
$97.7\%$ for NVDA. These findings supported the HMM variants as sources of heavy-tailed
asset paths whose holdout fit varied by asset. On the SPY holdout,
their advantage over the Laplace baseline lay in kurtosis rather than
in pass rates or distances. The jump mechanism lowered autocorrelation error mainly in
training and at a small cost in pass rates. The six-method composition comparison used the no-jump variant to
keep per-ticker jump calibration out of the comparison, and a separate
multi-asset comparison evaluated the transfer of jump settings.

% Main figure fig:model_comparison is collected after the references.

\subsection{Variance-corrected multi-asset composition}

We compared six composition methods on the $423$ non-market assets of
the training universe to test whether hybrid composition preserved
generator variance while recovering the calibrated market loading
(Table~\ref{tab:aggregate}; Fig.~\ref{fig:preservation}). The
GARCH(1,1)-$t$ method covered the $393$ assets with eligible fits.
With observed training SPY as the market path, the index trackers QQQ
and SPYG used the $R^2$-preserving branch, and the other $421$ assets,
which we call non-tracker assets, used the variance-preserving branch
without clipping (Table~\ref{tab:branch}). The two trackers targeted
the calibrated $R^2$ rather than the generator variance, so we marked
them separately (Fig.~\ref{fig:preservation}). We paired each composed
path with its uncomposed generator draw to measure the ratio of
composed to generator variance directly. Across the non-tracker
assets, the median of the per-ticker median ratios was $1.3171$ for
naive composition and $1.0001$ for hybrid composition, and the hybrid
ticker medians ranged from $0.9933$ to $1.0072$
(Fig.~\ref{fig:preservation}b). Individual hybrid paths still varied
around the target, because the correction assumes zero covariance
between the scaled asset draw and the market path, whereas a finite
path has a small nonzero sample covariance. Across non-tracker hybrid
paths, the median absolute departure of the variance ratio from one
was $1.12\%$, and its $95$th percentile was $3.41\%$
(Fig.~\ref{fig:preservation}b). The correlation between generator
draws and the observed market path averaged near zero for every asset
(Fig.~\ref{fig:cov_diag}), which supported the zero-covariance
assumption on average but not on each path.

% Main figure fig:preservation is collected after the references.

% Main table tab:aggregate is collected after the references.

The methods differed more in distributional fit than in factor
recovery (Table~\ref{tab:aggregate}; Fig.~\ref{fig:tails}). All six
methods had median absolute $\beta$ errors between $0.017$ and
$0.022$. The naive and hybrid methods had a median absolute intercept
error of $0.002\,\mathrm{yr}^{-1}$, but this small error followed from
centering, not from any ability of the generators to reproduce each
asset's mean. The OLS fit to observed data passes through the sample
means, so without clipping, the centered composed paths on the
observed training SPY path reproduced each observed asset mean
exactly. Refitting OLS to a composed path then gave an intercept error
equal to the loading error times the market mean. The three
residual-fit methods (JumpHMM-on-residuals, block bootstrap, and
GARCH(1,1)-$t$) resampled or fitted the OLS residual series. These
methods and Gaussian SIM did not center their draws, and their
intercept errors were an order of magnitude larger. Gaussian SIM had
the lowest median $R^2$ error but almost never passed the distribution
tests, and its median excess kurtosis was an order of magnitude below
the observed median of about $13$ across the $423$ assets. Hybrid
composition passed $73.6\%$ of KS comparisons, against $7.1\%$ for
naive composition, and its advantage held across the calibrated
$\beta$ range. Its own pass rate nonetheless fell from $87.6\%$ in the
lowest $\beta$ quartile to $65.9\%$ in the highest
(Fig.~\ref{fig:preservation}a; Table~\ref{tab:beta_bucket}). It retained heavy tails,
with median excess kurtosis of $7.2$ and a median Hill estimate of
$0.37$, reported as the reciprocal tail index so that larger values
indicate heavier tails (Fig.~\ref{fig:tails}). Its kurtosis was
similar to that of naive composition, so the variance rescaling did
not cause its shortfall from the observed median. Its median $R^2$ error fell between those of the naive and
Gaussian methods, because the variance-preserving branch targets the
generator variance rather than the observed $R^2$.

Fitting the OLS residuals directly needs no correction, because those
residuals already exclude the market component, and two of the three
residual-fit methods were competitive with hybrid composition
(Tables~\ref{tab:aggregate} and~\ref{tab:centering_control}). The
six-method comparison used each method's native mean treatment: naive
and hybrid paths were centered to a zero residual sum, whereas the
other four methods kept each draw's realized residual mean. Under
native treatment, the residual-fit KS pass rates ranged from $67.4\%$
for GARCH(1,1)-$t$ on its $393$ assets to $75.8\%$ for
JumpHMM-on-residuals. The hybrid rate lay about two percentage points
below JumpHMM-on-residuals and slightly above block bootstrap
(Table~\ref{tab:aggregate}). The JumpHMM-on-residuals and
block-bootstrap methods had slightly lower Wasserstein distances than
hybrid composition, and JumpHMM-on-residuals matched the hybrid
kurtosis and Hill estimate almost exactly. The block-bootstrap and
GARCH residual fits produced lighter tails than hybrid composition and
JumpHMM-on-residuals, with lower kurtosis and Hill estimates.
Subtracting each residual mean over the full training horizon from all
six methods, on the common $393$-asset set, changed the KS ranking
(Table~\ref{tab:centering_control}). On that set, the hybrid KS pass
rate was $74.4\%$ with or without centering. The JumpHMM-on-residuals
rate, $75.8\%$ on this set as on the full set, rose to $78.4\%$, and
the block-bootstrap rate rose from $73.5\%$ to $76.7\%$. Those two
methods also had higher AD pass rates and lower Wasserstein distances
than hybrid composition after centering, while Gaussian SIM and GARCH
remained below it on both pass rates. Centering brought every method's
median intercept error to $0.002\,\mathrm{yr}^{-1}$ at the reported
precision and left each path's variance, recovered loading, $R^2$, and
excess kurtosis unchanged to numerical tolerance. The pass-rate gains
therefore came from removing each draw's realized residual mean. Hybrid
composition kept its advantage over naive composition under this
control because both were already centered. The correction therefore improved on naive composition without a
second generator fit, but hybrid composition trailed
JumpHMM-on-residuals under both mean treatments and also trailed block
bootstrap after centering.

The emission fallback rules were triggered only in full-return
generators, so they could have handicapped naive and hybrid
composition relative to the residual-fit methods
(Tables~\ref{tab:fallback_assets} and~\ref{tab:fallback_sensitivity}).
Each of $39$ non-market models had one state filled by days with
unchanged closing prices, whose zero growth rates gave the state a sample spread below the fallback threshold. The
fallback replaced that spread with the asset's full-series scale. The same $39$ models held
$11$ further fallback states that were empty or held one observation,
and these kept the fallback throughout. The SPY model and every
JumpHMM-on-residuals fit needed no fallback, and the block-bootstrap
and GARCH methods have no emission states. To measure the effect, we gave each zero-return state its own near-zero empirical spread,
which stopped that state from emitting widely spread draws around
zero. We then replayed
the same $100$ training paths per affected asset with all other model
components and random draws unchanged. The change raised hybrid KS
pass rates from $68.6\%$ to $78.9\%$ and AD pass rates from $56.2\%$
to $70.2\%$ on those assets (Table~\ref{tab:fallback_sensitivity}).
The fallback rule therefore lowered the training fit of the
full-return methods on these assets without affecting the residual-fit
comparators. Spread over all $423$ assets, the change would have
raised the aggregate hybrid KS pass rate by about one percentage
point, less than its gap to JumpHMM-on-residuals. The main comparisons
retain the original rule, because the alternative was a post hoc
training diagnostic that was not evaluated on held-out histories. Both the main comparison and this check used training
histories only; the holdout evaluation tested whether the advantage of
hybrid over naive composition persisted beyond the fitting period.

\subsection{Holdout performance and VaR coverage}

To test out-of-sample generalization without refitting or look-ahead
bias, we froze the fitted models and SIM parameters before evaluating
them against $416$ complete asset histories from $2025$
(Table~\ref{tab:oos_composition}). Hybrid composition reached an
$87.0\%$ mean cross-ticker KS pass rate, against $83.5\%$ for naive
composition. Its median ratio of synthetic to observed $2025$ variance
was $0.97$, whereas naive composition inflated that ratio to $1.29$.
The JumpHMM-on-residuals method reached an $86.2\%$ KS pass rate and a slightly lower median $W_1$ than hybrid composition. As in
training, these holdout comparisons used each method's native mean
treatment, and no common-centering control was run on the holdout. Every method
passed more often in $2025$ than in the full-length training
comparison, and the gap between hybrid and naive composition shrank
(Tables~\ref{tab:aggregate} and~\ref{tab:oos_composition}). For every method, the pass rate against matched-length $249$-day
training blocks lay between its full-length training rate and its
$2025$ rate. The shorter sample, which lowers test power, therefore
accounted for much of the shrinking gap and part of the rise, although
the matched blocks were also scored with simulated rather than
observed SPY. The rest of the rise in $2025$ pass rates remained
unexplained. The average pass rates also concealed substantial
differences among tickers (Fig.~\ref{fig:oos_composition}). Most
ticker-level rates were high, but every method had a lower tail of
weakly fitted tickers. For hybrid composition, about three quarters of
tickers passed more often in $2025$ than against their matched-length
training blocks, and about one quarter passed less often
(Fig.~\ref{fig:oos_composition}b). The rise from training blocks to
$2025$ therefore did not hold for every asset.

To test whether the holdout improvement over naive composition
extended to reliable tail thresholds, we estimated per-ticker one-day
left-tail Value-at-Risk (VaR) thresholds from $5{,}000$ separate
$249$-day paths and counted exceedances on the observed $2025$
histories (Table~\ref{tab:var}). Doubling the threshold pool from
$2{,}500$ to $5{,}000$ paths changed every method's mean exceedance
rate by less than $0.006$ percentage points at either coverage level
(Table~\ref{tab:var_convergence}). The coverage level
$\alpha=0.95$ has a nominal exceedance rate of $1-\alpha=5\%$ and gave
each asset about $12$ expected exceedances. At that level, hybrid
composition had a mean
exceedance rate of $5.48\%$, closer to nominal than the $4.04\%$ of
naive composition. At $\alpha=0.99$, the hybrid rate of $1.12\%$ and
the naive rate of $0.88\%$ lay equally close to the $1\%$ nominal rate
on opposite sides. The JumpHMM-on-residuals method had rates of
$5.62\%$ and $1.16\%$, close to hybrid composition, and Gaussian SIM
had the largest $99\%$ rate at $1.72\%$. The advantage of hybrid over
naive composition thus appeared at the $95\%$ level but not at the
$99\%$ level. Neither comparison established a ranking, because each
asset contributed few exceedances, only about $2$--$3$ expected at the
$99\%$ level, and all assets shared one market period, so their
exceedances were not independent. The hybrid mean rates exceeded
nominal at both coverage levels, so its thresholds were slightly too
permissive on average, and the descriptive check could not confirm
their calibration. The
cross-asset standard deviations of the exceedance rates described how
realized rates varied among assets; they did not test miscalibration
or validate coverage in other periods.

% Main table tab:oos_composition is collected after the references.

% Main figure fig:oos_composition is collected after the references.

% Main table tab:var is collected after the references.

\subsection{Jump-enabled multi-asset composition}
\label{sec:results-jump-composition}

To test whether jump episodes improved the temporal fit of the
composed paths, we compared jumps
disabled, market-only jumps, and market-plus-asset jumps under naive
and hybrid composition (Table~\ref{tab:jump_ablation}, where hybrid
composition is labeled Corrected). Enabled generators used the SPY
jump settings without per-asset tuning. This comparison simulated the market path in both windows and used
each simulated path's variance in the correction, whereas the
six-method comparisons used observed SPY in training and the training
market variance in the holdout. With these conventions and different
seeds and replication counts, its no-jump pass rates differ by up to
about one and a half percentage points from those of the six-method
comparisons. Its temporal error is the path-averaged
ACF-MAE, which differs from the ensemble-average error of the
single-asset comparison, so the two are not directly comparable. With
hybrid composition, market-only jumps reduced the training ACF-MAE
over lags $1$--$25$ by $5.3\%$, and the $95\%$ Monte Carlo interval for that reduction excluded zero.
The KS pass rate fell by only $0.3$ percentage points, and the
interval for that change included zero
(Table~\ref{tab:jump_ablation_uncertainty}). Enabling
jumps in both the market and asset generators reduced the ACF-MAE by
$14.6\%$ from the no-jump setting but lowered the KS pass rate from $73.9\%$ to
$61.7\%$. That cost was far larger than the one jumps imposed on SPY,
whose settings were tuned. Across all settings and both windows, the mean
ratio of composed to generator variance under hybrid composition
differed from one by less than $0.05\%$. The correction therefore held
the variance target with jumps enabled; both jump settings lowered
training autocorrelation error, but only adding jumps to the asset
generators materially lowered the KS pass rate. The training-period
temporal improvement did not carry into the $2025$ holdout
(Table~\ref{tab:jump_ablation}). With hybrid composition, market-only
jumps raised the $25$-lag ACF-MAE by $2.7\%$ and market-plus-asset
jumps raised it by $6.9\%$. The $95\%$ Monte Carlo intervals excluded
zero for both increases, and the $60$-lag errors also rose
(Tables~\ref{tab:jump_ablation}
and~\ref{tab:jump_ablation_uncertainty}). The holdout KS pass rate
changed little under market-only jumps and fell by about $1.5$
percentage points under market-plus-asset jumps. Only $2.2\%$ of the
simulated one-year market paths contained an episode
(Table~\ref{tab:jump_ablation_episodes}). The holdout therefore judged
the transferred settings on few episode paths and a single year of
observed data, and it could not establish their behavior in other
market periods.

% Main table tab:jump_ablation is collected after the references.

\subsection{Sensitivity analyses}

We examined how the branch-selection threshold
$R^2_{\mathrm{preserve}}$ assigned assets to branches in the observed
universe and whether any value of the variance floor $f$ triggered
clipping (Figs.~\ref{fig:r2_distribution}
and~\ref{fig:branch_map}; Table~\ref{tab:branch}). The
$R^2$-preserving branch applied only to QQQ and SPYG, whose hybrid
compositions passed every KS test in training. The highest-$R^2$ asset
outside these trackers, the individual stock BLK, fell about $0.16$
below the $0.80$ threshold, so the assignment did not depend on a
narrow margin. Across the explored thresholds from $0.70$ to $0.90$ at
the default floor, only the QQQ assignment changed, moving to the
variance-preserving branch at $0.90$. At the default threshold, no
value of $f$ from $0.05$ to $0.30$ caused any asset to be clipped. The
empirical evaluation of the $R^2$-preserving branch was therefore
limited to two assets, and the observed market variance never
activated clipping.

Because only two real assets used the $R^2$-preserving branch, we also
checked it on a synthetic grid with known loading and $R^2$
(Table~\ref{tab:tracker_grid}). The grid covered
$(\beta_{\mathrm{true}},R^2_{\mathrm{true}})\in\{0.8,1.0,1.2\}
\times\{0.80,0.85,0.90,0.95,0.99\}$, with $100$ replications per cell.
In each replication, a synthetic tracker series equaled
$\beta_{\mathrm{true}}$ times the observed training SPY path plus an
independent Gaussian residual, with zero intercept and the residual
variance that gives $R^2_{\mathrm{true}}$ under
Eq.~\eqref{eq:r2-target}. The composition algorithm received
$\beta_{\mathrm{true}}$ and $R^2_{\mathrm{true}}$ as its calibrated
inputs and a second, independent Gaussian draw in place of a generator
draw. The composer's generator variance was fixed at the target
residual variance rather than recomputed from the draw, so the
residual scale in Eq.~\eqref{eq:r2-rescale} was exactly one. Median
loading and $R^2$ estimates from the composed series were close to
their targets. Every composed series also passed the KS comparison
with its synthetic tracker. Both series contained the same SPY path,
however, so the two samples were not independent and these pass rates
were descriptive rather than calibrated tests. The grid therefore
checked the implementation and finite-sample accuracy of the
$R^2$-preserving formula under a restricted Gaussian construction; it
did not exercise full-return HMM fitting, non-unit residual scaling,
or simulated market paths.

The decline of hybrid fit with calibrated $\beta$
(Table~\ref{tab:beta_bucket}; Fig.~\ref{fig:preservation}a) follows
the market loading ratio $\rho_i$, the market's share of the target
variance, which a larger loading raises for an asset of given
variance. The scaled draw is a full-return path that already carries
the asset's whole shape, and composition adds an independent market
term to it. As $\rho_i$ grows, the market's shape therefore makes up
more of the blend, and the composed distribution departs further from
the asset's observed one; Online
Appendix~\ref{sec:supp-residual-choice} derives this dilution for
kurtosis. A more volatile market raises $\rho_i$ in the same way and,
once $\rho_i$ exceeds $1-f$, activates clipping, which the observed
market variance never did. To measure when clipping activates, we
multiplied deviations of the observed training SPY path from its mean
by $\gamma\in\{1,2,3\}$, which preserved the market mean and
multiplied its standard deviation by $\gamma$. We then recomposed all
$423$ assets with the scaled path's variance in the correction,
holding the fitted marginals, intercepts, and calibrated loadings
fixed (Table~\ref{tab:stress}). Clipping applied to no composed path
at $\gamma=1$ and to $76.4\%$ and $95.2\%$ of composed paths at
$\gamma=2$ and $3$. Because clipped paths retained their fitted
intercepts while the market mean stayed unchanged and nonzero, their
means shifted by $(\beta_i^{\eff}-\beta_i)\overline{g}_m$ by
construction. The stress test therefore characterized when clipping
activated and how it moved asset means; it did not evaluate fit under
the scaled markets.
